% !TEX root = ../main.tex
% !TEX program = pdflatex
\subsection{Определение степени ассоциативности}

В $k$-канальном наборно-ассоциативном кэше одна строка памяти может занять один из $k$ каналов своего набора. Чтобы проверить ассоциативность L1, программа строит цикл из $N$ фрагментов, начала которых разнесены на размер кэша $C$. Из каждого фрагмента используются два адреса внутри одной строки, поэтому в тесте участвуют ровно $N$ конфликтующих строк, но цепочка не вырождается в повторное чтение единственного адреса. Число фрагментов меняется от 1 до 32.

Для L1 все адреса с одинаковыми битами смещения в странице относятся к одному виртуально индексируемому набору. При шаге $C=128$~КБ страницы 16~КБ гарантируют совпадение индексных битов. Для L2 тот же тест не гарантирует попадание в один физический набор: физические номера страниц скрыты от пользовательской программы. Поэтому результат для L2 используется только как демонстрация ограничения метода.

\textbf{Кэш L1} ($C=128$~КБ). Результаты показаны на рисунке~\ref{fig:assoc_l1}. Для 4--8 конфликтующих строк медианное время составляет примерно 1,15--1,21~нс на переход. При $N=9$ оно возрастает до 2,23~нс, затем продолжает расти по мере увеличения доли промахов. Первая строка сверх восьми вызывает измеримый рост задержки, что согласуется с \textbf{8-канальной ассоциативностью L1}. Рост не обязан сразу достигать полной задержки L2: при $N=9$ часть обращений ещё попадает в кэш.

\fig{assoc_l1.png}{Конфликтный тест ассоциативности L1: минимум и среднее время}{fig:assoc_l1}

\textbf{Кэш L2} ($C=16$~МБ). Результаты показаны на рисунке~\ref{fig:assoc_l2}. Время постепенно растёт с увеличением числа фрагментов, но ясного порога $k+1$ нет. При этом одинаковые младшие биты виртуальных адресов также создают конфликты в L1, поэтому рост около $N=9$ нельзя приписать L2. Для L2 остаётся неизвестным, в какие физические наборы попали узлы: пользовательская программа не может выбрать физические номера страниц. Из-за этого кривая включает конфликты физического кэша, задержки L1 и возможные эффекты TLB, а не изолированную ассоциативность L2. Запрос больших страниц сам по себе не гарантирует нужного распределения по наборам.

\fig{assoc_l2.png}{Конфликтный тест с шагом 16 МиБ: результат не определяет ассоциативность L2}{fig:assoc_l2}

Таким образом, ассоциативность L1 оценивается как 8 каналов по выраженному росту после восьми конфликтующих строк. Ассоциативность L2 этим пользовательским тестом достоверно определить нельзя; прежнюю оценку «не менее 12, вероятно 16» по этому эксперименту не используем.
